\section{Whittaker values at exceptional finite places}
\label{sec:ramified-whittaker}

There are two local normalizations in the theta comparison. At a
nonsingular Hermitian frequency the Whittaker coefficient measures a
Gram fiber. In the split rank descent a one-dimensional Fourier
transform introduces a ramified Gauss sum. We compute both
normalizations. This distinguishes the scalar in a rank step from
the measure used to count solutions to a Gram equation.

\subsection{The weighted Gram fiber}

Let \(E/F\) be the CM extension with conjugation \(\sigma\), \(v\) any
finite place of \(F\), \(k=F_v\), and \(A=k\otimes_F E\). Use the
involution induced by \(\sigma\) and write \(b^*={}^{t}\bar b\).
Let \(J\) be a nonsingular Hermitian form in \(M\) variables over
\(E\). Let \(V=(A^M)^n\), \(H=\operatorname{Herm}_n(A)\), and let
\(Q(Y)\) be the Hermitian Gram matrix of the \(n\) vectors in \(Y\)
with respect to \(J\). Fix \(T\in H\) with \(\det T\in A^\times\).
This determinant condition also makes sense when \(A=k\times k\).
Write \(\psi_v\) for the local component of the specified nontrivial
unitary adelic additive character. On \(H\) the pairing used with
this character is
\[
 [b,D]=\tfrac12\operatorname{Tr}_{A/k}\operatorname{tr}(bD).
\]
Take the prescribed self-dual Haar measures \(\mu\) on \(V\) and
\(\eta\) on \(H\), with \(\eta\) self-dual for \(\psi_v([b,D])\).
Take any compact open additive subgroups \(L_1,L_2\subset A^n\) and
an additive Haar measure \(\lambda\) there. For \(b\in H\) set
\[
 h(b)=\frac{\lambda(L_1+bL_2)}{\lambda(L_1)},\qquad
 s_0=(M-n)/2 .
\]
Let \(\phi\) be locally constant and compactly supported on \(V\).
The density input is a complex-valued function \(g\), continuous on
the open set \(H^\times\) of nonsingular forms, such that
\[
 \int_{Q^{-1}(C)}\phi(Y)\,d\mu(Y)=\int_Cg(D)\,d\eta(D)
 \quad(C\subset H^\times\text{ compact}).
\]
The low-rank and equal-rank density theorems used by the global
formula supply such functions. The following local argument applies
to every \(g\) satisfying this condition.

Put
\[
 U_T(b)=\int_V\phi(Y)\psi_v([b,Q(Y)-T])\,d\mu(Y),
 \qquad F_z(b)=h(b)^{-(z-s_0)}U_T(b).
\]
There is an entire function \(W_T(z)\) equal to
\(\int_H F_z(b)\,d\eta(b)\) on a right half-plane, and
\begin{equation}
                       W_T(s_0)=g(T).
 \label{eq:exceptional-whit-density}
\end{equation}
We prove this also for nonstandard test functions and lattices.

\paragraph{Extraction of the density.}
Give matrices the maximum of their entry coordinate norms in a fixed
\(k\)-basis of \(A\). Write \(B_R=\{b\in H:\|b\|\leq R\}\) and
\(P_R=\{D\in H:\psi_v([b,D])=1\ \forall b\in B_R\}\). For \(R>0\),
\(B_R\) is a compact open additive subgroup. Orthogonality gives
\[
  \int_{B_R}\psi_v([b,D])\,d\eta(b)=\eta(B_R){\bf1}_{P_R}(D).
\]
Indeed translation by any \(b_0\in B_R\) multiplies this integral by
\(\psi_v([b_0,D])\). If \(D\notin P_R\), choose \(b_0\) with
nontrivial multiplier and subtract.

Here is an explicit bound on the annihilator. Choose \(y_0\in k\)
with \(\psi_v(y_0)\ne1\). Nondegeneracy of the trace form on the
separable algebra \(A\) gives \(c>0\) such that for each \(t\in A\)
some \(a\) satisfies \(\|a\|\leq1\) and
\(\lvert\operatorname{Tr}_{A/k}(at)\rvert\geq c\|t\|\).
For example, use a fixed trace-dual basis, scaled into the unit
ball by one scalar, and the equivalence of coordinate norms.
Choose \(C_\sigma\) with \(\|\bar a\|\leq C_\sigma\|a\|\).
For \(D\ne0\) choose a largest entry \(D_{ij}\) and apply the trace
bound to get \(a\) and \(w=\operatorname{Tr}(aD_{ij})\ne0\).
With \(X\) supported at \((j,i)\) with entry \(a\), the matrix
\(b=(y_0/w)(X+X^*)\) is Hermitian. Trace invariance under the
involution and cyclicity give \([b,D]=y_0\). Thus
\[
 D\in P_R\quad\Longrightarrow\quad
 \|D\|\leq\frac{|y_0|\max(1,C_\sigma)}{cR}.
\]
The factor \(1/2\) in the pairing cancels the two traces of \(X\)
and \(X^*\), including in residue characteristic two.
In particular \(P_R\) is closed and bounded, hence compact.
Fourier inversion at zero applied to \({\bf1}_{B_R}\), whose transform
is the displayed compactly supported function, now gives
\(\eta(B_R)\eta(P_R)=1\).

Fubini applies to \(B_R\times\operatorname{supp}\phi\). It gives
\[
 \int_{B_R}U_T(b)\,d\eta(b)=\eta(B_R)
                  \int_{Q^{-1}(T+P_R)}\phi(Y)\,d\mu(Y).
\]
For all sufficiently large \(R\), the compact set \(T+P_R\) lies in
\(H^\times\), and the last integral equals \(\int_{T+P_R}g\,d\eta\).
The absolute difference between the left side and \(g(T)\) is at
most \(\sup_{D\in T+P_R}|g(D)-g(T)|\), by the volume product above.
Continuity at \(T\) makes this tend to zero.

\paragraph{Cancellation on a large shell.}
Fix \(q_0>1\). We show that for every sufficiently large \(R\) and
every complex \(z\),
\[
          \int_{R<\|b\|\leq q_0R}F_z(b)\,d\eta(b)=0 .
\]
For \(a=1+x\) with \(\|x\|\leq r_*\) sufficiently small, the
transform \(b\mapsto a^*ba\) preserves each norm ball and its Haar
measure. Indeed choose \(r_*\) so the correction
\(x^*b+bx+x^*bx\) has norm strictly less than \(\|b\|\) for \(b\ne0\),
and \(1+x\) has an inverse. Preservation of a compact open ball of
positive volume fixes the Haar scaling at one. Also \(aL_2=L_2\)
and \(a^*L_1=L_1\), since both lattices are compact open and these
automorphisms and their inverses approach the identity. It follows
that \(h(a^*ba)=h(b)\). The integral
\(J_\phi(b)=\int_V\phi(Y)\psi_v([b,Q(Y)])\,d\mu(Y)\) is unchanged:
make the corresponding change of basis on the tuple \(Y\). A single
congruence neighborhood fixes \(\phi\) under this change. One can
choose it by expressing \(\phi\) as a finite linear combination of
indicators of cosets of one compact open subgroup and stabilizing
those cosets. Shrinking \(r_*\) ensures all of these properties.

The only varying factor of \(F_z\) is now the frequency character.
Expansion gives
\[
 [T,a^*ba]-[T,b]
 =\operatorname{Tr}_{A/k}\operatorname{tr}(Tbx)+[T,x^*bx].
\]
The two terms linear in \(x,x^*\) have the same field trace, which
cancels the \(1/2\) in the pairing. There exists \(C_Q\geq0\) such
that \(\lvert[T,x^*bx]\rvert\leq C_Q\|x\|^2\|b\|\), by continuity
and finite dimensionality of multiplication, trace and involution.
Choose \(\varepsilon>0\) with \(\psi_v\) trivial on
\(\{|t|\leq\varepsilon\}\).

The Fourier character in the linear term can be made nontrivial
uniformly in the direction of \(b\). Choose \(C_m\geq0\) with
\(\|AB'\|\leq C_m\|A\|\|B'\|\) for matrices and set
\[
 C_\nu=\frac{|y_0|(C_m\|T^{-1}\|+1)}{c}.
\]
For \(b\ne0\), let \(Y=Tb\), choose a largest entry \(Y_{ij}\) and
choose \(a'\) of norm at most one with
\(\lvert w\rvert=|\operatorname{Tr}(a'Y_{ij})|\geq c\|Y\|\).
Take \(x_0\) supported at \((j,i)\) with entry \(-y_0a'/w\).
Then
\[
 \psi_v(-\operatorname{Tr}\operatorname{tr}(Tbx_0))=\psi_v(y_0)\ne1,
 \qquad \|x_0\|\leq C_\nu/\|b\|.
\]
The bound follows by writing \(b=T^{-1}Y\).

Take
\[
 R\geq\max\{1,C_\nu/r_*, C_Q C_\nu^2 q_0/\varepsilon\},
 \qquad r=C_\nu/R .
\]
On the shell and \(\|x\|\leq r\) the quadratic term lies in the
kernel of \(\psi_v\). Consequently substitution by \(1+x\) multiplies
the integrand at \(b\) by precisely
\(\psi_v(-\operatorname{Tr}\operatorname{tr}(Tbx))\).
Average the substitutions over the additive matrix ball of radius
\(r\), with normalized Haar measure. Each substitution preserves the
shell integral. For each \(b\) in the shell the multiplier has mean
zero, by translation by the \(x_0\) just constructed. Fubini on the
compact sets proves the cancellation.

\paragraph{The half-plane and continuation.}
Choose \(0<|\pi|<1\) in \(k\), put \(q_1=|\pi|^{-1}\), and let
\(K>0\) be the Haar scaling \(\eta(B_{rq_1})=K\eta(B_r)\) of
multiplication by \(\pi^{-1}\). There is \(R_1>0\) such that
\[
       \|b\|\geq R_1q_1^j\quad\Longrightarrow\quad h(b)\geq 2^j
       \qquad(j\geq0).
\]
Indeed take \(\rho\geq0\) bounding the norm on \(L_1\) and take a
small nonzero scalar \(t_0\) with \(u t_0 e_i\in L_2\) for each
coordinate vector \(e_i\) and every \(u\in\mathcal O_k\). Put
\(R_1=(\rho+1)/|t_0|\). Choose a column \(b e_i\) containing an
entry of norm \(\|b\|\). For every binary word of length \(j\) take
\(u=\sum_{\ell=0}^{j-1}d_\ell\pi^\ell\), \(d_\ell\in\{0,1\}\).
Distinct words differ in norm by at least \(|\pi|^j\), by their
first differing digit. The \(2^j\) points \(u t_0 b e_i\) are
distinct modulo \(L_1\), since their differences have norm at
least \(|\pi|^j|t_0|\|b\|\geq \rho+1\). Their \(L_1\)-cosets
all lie in \(L_1+bL_2\), which proves the asserted lower bound.

Write \(s=\Re z-s_0\) and \(V_1=\eta(B_{R_1})\).
For \(s>\log_2(2K+2)\) we have \(K2^{-s}<1\). On the shell
\(R_1q_1^j\leq\|b\|\leq R_1q_1^{j+1}\), the integral
of \(|F_z|\) is at most
\(\|\phi\|_{L^1(\mu)}K V_1(K2^{-s})^j\). The compact ball adds a
finite contribution. This proves absolute convergence on the
stated half-plane.

Apply shell cancellation with \(q_0=2\). Fix \(R\) beyond its cutoff.
For every \(z\) the integrals over \(B_{2^jR}\) all equal the integral
over \(B_R\). On this compact ball \(1\leq h\leq C\), since \(L_2\)
and the ball are compact and the Minkowski sums lie in a bounded
set. The index \(h\) is measurable, indeed locally constant, by
local constancy of the sets \(L_1+bL_2\) in \(b\).
Differentiation under the compact integral is valid uniformly on
compact \(z\)-sets, since each derivative contributes a bounded
power of \(\log h\). The integral over \(B_R\) therefore defines
an entire function. On the half-plane of absolute convergence the
exhaustion gives the integral over \(H\). At \(s_0\) the weight is
one, so the exhaustion and the density limit give
\eqref{eq:exceptional-whit-density}. The local computation applies
at each member of the finite exceptional set in the global
factorization, independently of any formula for spherical vectors.

\subsection{The conductor and phase in split rank descent}

Now take \(v\in T=S_0\cup Q\) as in Section~\ref{sec:theta-match}.
This is a split place with designated embedding \(E\to k\). Write
\(\nu\) for the corresponding local component of the ratio of the
two characters. It is unitary and ramified, and \(\psi_v\) has
annihilator \(\mathcal O_k\) on \(\mathcal O_k\). Write \(\varpi\)
for a uniformizer, \(q=|\mathcal O_k/\varpi|\), \(U_0=\mathcal
O_k^\times\), and \(U_i=1+\varpi^i\mathcal O_k\) for \(i>0\).
Let \(f\) be the least \(i\) with \(\nu(U_i)=1\). It is positive.
Continuity gives a finite conductor: choose \(U_i\) on which all
values lie within \(1/2\) of one. Every power of each of its
character values then lies in the same disc. A subgroup of the
circle contained in this disc is trivial. Minimality gives a
\(u_0\in U_{f-1}\) with \(\nu(u_0)\ne1\), also when \(f=1\).

Put \(D_f=\mathcal O_k/\varpi^f\mathcal O_k\). For \(x\in D_f\) define
\[
 G(x)=\sum_{y\in D_f^\times}\nu(\widetilde y)^{-1}
                     \psi_v(\widetilde x\widetilde y/\varpi^f).
\]
The tilde denotes any integral lift. All lifts of a unit are units,
and replacing one by another multiplies it by an element of \(U_f\).
Changes of either lift change the additive argument by an integer
of \(k\). Thus both factors are well defined. In additive Haar
measure \(dt\) with \(\operatorname{vol}(\mathcal O_k)=1\), one has
the exact integral
\[
 \int_{\mathcal O_k^\times}\nu(t)^{-1}
                  \psi_v(\widetilde x t/\varpi^f)\,dt=q^{-f}G(x),
\]
since each additive coset in the support has mass \(q^{-f}\).

For a unit \(a\), the substitution \(y\mapsto ay\) gives
\(G(ax)=\nu(a)G(x)\). For a nonunit \(x\), multiplication of \(x\)
by \(u_0\) does not change \(x\bmod\varpi^f\), so
\(G(x)=\nu(u_0)G(x)\), hence \(G(x)=0\).
The additive pairing on \(D_f\) is perfect by the conductor
assumption on \(\psi_v\): if
\(\psi_v(xd/\varpi^f)=1\) for every integral \(x\), then
\(d/\varpi^f\in\mathcal O_k\). Translation on the finite group
therefore gives
\[
 \sum_{x\in D_f}\psi_v(xd/\varpi^f)
     =\begin{cases}q^f&d=0\text{ in }D_f,\\0&d\ne0\text{ in }D_f.
       \end{cases}
\]
Multiply \(G(x)\) by \(\psi_v(-x/\varpi^f)\) and sum in \(x\).
Only the term \(y=1\) survives, and gives \(q^f\). This proves the
source's nonvanishing \(\operatorname{gauss}(\psi_v,\nu)=G(1)\ne0\).

The same finite calculation gives the size when the characters are
unitary. Expand \(\sum_x|G(x)|^2\) and sum in \(x\) first.
Orthogonality leaves \(q^f|D_f^\times|\). Summing instead over
units \(x\) and using \(G(x)=\nu(x)G(1)\) gives
\(|D_f^\times|\,|G(1)|^2\). Thus
\[
 |G(1)|=q^{f/2},\qquad
 w_v=\nu(\varpi)^f\,q^{-f/2}G(1),\qquad |w_v|=1 .
\]
At the integer argument \(z\) of the source's rank step, its local
factor with the chosen sign \(e\in\{1,-1\}\) is exactly
\begin{equation}
 \nu(e)\epsilon_v(z)=\nu(e)\,q^{f(1/2-z)}w_v .
 \label{eq:rankstep-local-phase}
\end{equation}
Indeed the definition of \(\epsilon_v(z)\) is
\(q^{-fz}G(1)\nu(\varpi)^f\). Replacing \(\varpi\) by \(a\varpi\)
for a unit \(a\) changes the sum by \(\nu(a)^{-f}\) and the
uniformizer character factor by \(\nu(a)^f\). These cancel.
Replacing \(\psi_v(t)\) by \(\psi_v(at)\) multiplies \(w_v\) by
\(\nu(a)\). The character \(\nu(e)\) squares to one.
The hypotheses give the phase \(w_v\) by the displayed sum. A
numerical root of unity for this phase would require further
specialization. Even a finite-order character on units does not
specify the value of \(\nu(\varpi)\).

For the beta and alpha steps the argument is respectively
\(z=-\delta\) and \(z=1-\delta\), where
\(\delta=(a-a')/2\). Each exceptional place contributes
\eqref{eq:rankstep-local-phase}, with the sign of its rank-step
formula. The finite character sum computes these character-dependent
factors. The geometric factor in the rank step, the surjectivity
of the two maps onto the split induced representation and the
identity of local tensor models are separate inputs proved in the
respective descent and local matching arguments.

\subsection{A matched pair and the sign of the cone integral}

We can determine the local sign in \eqref{eq:rankstep-local-phase}
from the standard vectors used in the source. Put
\(\ell_\nu(t)={\bf1}_{\mathcal O_k^\times}(t)\nu(t)^{-1}\), extended
by zero at \(t=0\). If \(e\) selects \(r\) columns of an \(r\)-by-\(d\)
matrix \(Y\) in their indicated order, set
\[
  S_e(Y)=\ell_\nu(\det Y_e)
  \prod_{i,j}{\bf1}_{L_j}(Y_{ij}),\qquad
  L_j=\begin{cases}
       \mathcal O_k&j\in e,\\
       \varpi^f\mathcal O_k&j\notin e .
       \end{cases}
\]
Let \(e_A\) select the first \(a\) columns and \(e_B\) the last \(a'\).
These standard vectors are compactly supported, locally constant,
and have equal nonzero images in the two transforms \(P_A,P_B\)
of Section~\ref{sec:theta-match}.

For a matrix \(y\) of size \(a'\)-by-\(a\), write
\(u(y)=\left(\begin{smallmatrix}1&0\\y&1\end{smallmatrix}\right)\).
The matrices inserted in the two integrals have frames
\([1,-{}^ty]\) and \([y,1]\). In either transform the selected
minor is \(g\). On the support of the integrand its entries are
integral and its determinant a unit, so \(g\in\mathrm{GL}_r(\mathcal
O_k)\). The inverse is integral too. Thus the off-minor condition
is equivalent to \(y\in\varpi^f M_{a',a}(\mathcal O_k)\). The
weight \(\nu(\det g)\) cancels \(\ell_\nu(\det g)\), both absolute
determinant factors are one and the compact group has measure one.
Both transforms therefore give the indicator of this same set in
\(y\), with value one at zero.

To verify equality at arbitrary \(h\), write
\(h=\left(\begin{smallmatrix}A&B\\C&D\end{smallmatrix}\right)\).
When \(D\) is singular, its complementary minor in \(h^{-1}\)
is singular as well, by the complementary-minor formula for an
invertible matrix. Thus each selected determinant vanishes in the
corresponding frame and both integrands are zero. When \(D\) is
invertible put \(y=D^{-1}C\). Then \(h=p u(y)\) with \(p\) upper
block triangular. In the two transforms, left multiplication by \(p\)
has the same scalar. Indeed, writing \(A_0,D_0\) for its diagonal
blocks, changes of the integration variables by \(A_0^{-t}\) on
the large side and by \(D_0\) on the small side give
\[
 \alpha(\det A_0)\beta(\det D_0)
            |\det A_0|^{a'/2}|\det D_0|^{-a/2}.
\]
Here \(\nu=\alpha/\beta\). The covariance on the large side uses
the ramified shell transformation law of \(P_A\).
This proves equality as functions on all of \(\mathrm{GL}_d(k)\).
Transport the vectors to the two Weil models and write them
\(\phi,\phi'\). The construction at the lower rank gives matched
vectors \(\phi_1,\phi'_1\) with common nonzero image.

In the beta step \(a'=b+2\), put \(r=a\) and \(j=\delta+1\).
In the alpha step \(a=b+2\), put \(r=a'\) and \(j=\delta\). The
zero-orbit evaluation of the standard vector is \(q^{-fr}\) times
its lower standard vector: its \(r\) independent restricted
coordinates range over \(\varpi^f\mathcal O_k\), and the remaining
integrated unit balls have mass one. We compute the cone evaluation
in the same coordinates. Number its distinguished rows and selected
columns from \(0\) to \(b+1\), put \(l=b+1\), and let \(e\) select
the block undergoing rank descent. In the beta chart the free
coordinate positions \(P\) are those in row \(l\), together with
those in column \(e(l)\) outside row \(0\). In the alpha chart
they are those in column \(e(0)\), together with those in row \(0\)
outside column \(e(l)\). The corner \((0,e(l))\) is fixed to the
first split variable \(s\), and the opposite corner \((l,e(0))\)
has an integration variable \(u\) with character \(\psi_v(-tu)\),
where \(t\) is the second split variable.

In the beta chart, expansion along row \(0\) and then the column
\(e(0)\), and in the alpha chart expansion along column \(e(l)\)
and then row \(l\), give the same selected determinant
\[
                          -su\det(Y_{\rm mid}).
\]
The sign is the sign of interchanging the first and last of the
selected columns. On the integral domain of the coordinate
indicators all factors of this product are integral, so the
determinant is a unit exactly when they all are. Its character
factor splits as
\(\ell_\nu(s)\ell_\nu(-u)\ell_\nu(\det Y_{\rm mid})\).
Change \(u\) to \(-u\) in the corner integral. The two minus signs
cancel, giving
\[
       I(t)=\int_{\mathcal O_k^\times}\nu(u)^{-1}\psi_v(tu)\,du .
\]
There are exactly \(r\) free positions with column outside \(e\):
the beta chart has \((l,e_A(i))\) for the \(a\) unselected columns, and the alpha chart
has \((0,e_B(i))\) for the \(a'\) unselected columns.
Each contributes \(q^{-f}\). All other coordinate integrals carry
unit volume. The surviving factor at the split variables is
therefore
\[
                       q^{-fr}\ell_\nu(s) I(t)
\]
times the lower standard vector.

For completeness, \(I(t)\) is supported on
\(\varpi^{-f}\mathcal O_k^\times\). For
\(|t\varpi^f|<1\) use \(u_0\in U_{f-1}\) as above: substitution
by \(u_0\) preserves \(\psi_v(tu)\) and multiplies the integral
by a nontrivial scalar. For \(|t\varpi^f|>1\) choose
\(h\in\varpi^f\mathcal O_k\) with \(\psi_v(th)\ne1\), using the
annihilator of \(\mathcal O_k\). Translation by \(h\) preserves
\(\ell_\nu\) and forces the integral to vanish. For a unit \(v'\)
one has \(I(\varpi^{-f}v')=\nu(v')I(\varpi^{-f})\), and
\(I(\varpi^{-f})=q^{-f}G(1)\) by the coset integral.

The two components of the split idele have character
\(\nu(s)\nu(t)^{-1}\) and norm \(|s|\,|t|\). The local integral
uses multiplicative Haar measure giving the product of the unit
groups mass one. Its support is a single translate of that subgroup,
\(s\in\mathcal O_k^\times\), \(t=\varpi^{-f}v'\) with \(v'\) a
unit. On it \(\ell_\nu(s)\nu(s)=1\), the two values of \(\nu(v')\)
cancel and \((|s||t|)^j=q^{fj}\). The local cone operator therefore
takes the standard vector to its lower vector times
\[
 q^{-fr}\, q^{-f}G(1)\nu(\varpi)^f q^{fj}
                  =q^{-fr}\epsilon_v(1-j).
\]
The compact stabilizer average has mass one and fixes the vector.
Comparison with the zero-orbit evaluation cancels exactly \(q^{-fr}\).
Writing \(p_1,p'_1\) for the lower quotient maps, the beta formula is
\[
 p'_1(\mathrm{cone}_v\phi')=\epsilon_v(-\delta)\,
                                p_1(\mathrm{ev}_v\phi).
\]
The alpha formula is
\(p_1(\mathrm{cone}_v\phi)=\epsilon_v(1-\delta)\,
p'_1(\mathrm{ev}'_v\phi')\). The respective vector on the right
before multiplication is nonzero because the lower standard vectors
match and the value of their common transform at the identity is one.

Thus the source's witnesses for the local sign \(e\) are \(+1\) in
both steps. To extend from these vectors to all matched fibers,
the source uses ramified multiplicity one for the equivariant
maps on the common quotient. In the beta case, let \(D,D'\) be
the linear maps it obtains by descending the evaluation and the
cone, respectively. On the common image \(\xi=p\phi=p'\phi'\)
the computation gives \(D'\xi=\epsilon_v(-\delta)D\xi\) and
\(D\xi\ne0\). Since \(D'=cD\) for some scalar, evaluation at
\(\xi\) forces \(c=\epsilon_v(-\delta)\). The same reasoning
with the roles interchanged gives \(\epsilon_v(1-\delta)\) in
the alpha case. Existence and equivariance of the descended maps,
the local split coordinate changes and ramified multiplicity one
are the representation-theoretic inputs to this last argument.

\paragraph{Correspondence with the release.}
In \texttt{SiegelWeilFormula.LocalCharacterIntegrals} the statements
used above are \texttt{shellP\_core}, \texttt{perp\_bound},
\texttt{vol\_mul\_vol\_perp}, \texttt{dualP}, \texttt{convP} and
\texttt{locFinX\_of}. The global applications are
\texttt{whitGenLtOf} and \texttt{whitGenEqOf} and import the
representation-density hypotheses. The split calculation uses \texttt{std\_pair}, \texttt{int\_cone}, \texttt{idele\_int}, \texttt{eps\_eq}, \texttt{stdB}, \texttt{stdA\_gen}, \texttt{locB\_of} and \texttt{locA\_of}. In \texttt{WeilTensorModels}
the functions are \texttt{gauss}, \texttt{eps} and \texttt{stepFac}.
\texttt{ThK.GaussNe.gauss\_ne\_zero} proves the required
nonvanishing and \texttt{ThKG.Loc.stepFac\_eq} identifies the finite
product. The modulus and the change-of-uniformizer computation above
follow by the finite sums displayed here.
